% !TEX root = ./Combinatorics Notes.tex
\documentclass[11pt]{scrartcl}

%%%%%%%%%%%%%% LATEX SAMPLE FILE %%%%%%%%%%%%%%%%
% A line which starts with a % sign
% is called a COMMENT. It is IGNORED
% by the LaTeX processor.

% Include math
\usepackage{amsmath,amssymb}
% Include links
\usepackage{hyperref}
\usepackage[dvipsnames]{xcolor}
\usepackage{tikz}
% \usetikzlibrary{positioning,chains,fit,shapes,calc}
\usepackage{forest}
% \usetikzlibrary{graphs,graphdrawing,arrows.meta}
\usepackage[utf8x]{inputenc}
\usepackage[dutch]{babel}
\usepackage{tikz}
\usetikzlibrary{shadows}
% \usegdlibrary{trees}
\usepackage{framed}
\usepackage[framed]{ntheorem}
\usepackage{xcolor}




%%%%%%%%%%%%%  THEOREMS  %%%%%%%%%%%%%%%%%
% Let's define some theorem environments
% To use later in the paper
\theoremstyle{plain} % other options: definition, remark

% By including [theorem], the lemma follows the numbering of theorem
% e.g. Thm 1, Lemma 2, Thm 3, Thm 4, \dots
\theoremstyle{definition}
% \newtheorem{definition}{Definition} % the star prevents numbering

% Remarks
\theoremstyle{remark}
\newtheorem{remark}{Remark}


\definecolor{formalshade}{rgb}{0.95,0.95,1}
\definecolor{darkblue}{rgb}{0,0,1}

% \newenvironment{proof}{%
%   \def\FrameCommand{%
%     \hspace{1pt}%
%     {\color{darkblue}\vrule width 2pt}%
%     {\color{formalshade}\vrule width 4pt}%

%   }%
%   \MakeFramed{\advance\hsize-\width\FrameRestore}%
%   \noindent\hspace{-4.55pt}% disable indenting first paragraph
%   \begin{adjustwidth}{}{7pt}%
%   \vspace{2pt}\vspace{2pt}%
% }
% {%
%   \vspace{2pt}\end{adjustwidth}\endMakeFramed%
% }

\theoremclass{Problem}
\theoremstyle{empty}
\theorembodyfont{\normalfont}
 \newcommand{\pmbox}[1]{
  \tikzstyle{pmbox} = [rectangle, rounded corners, draw=black, fill=green!5, inner sep=10pt, drop shadow={fill=black, opacity=.2}]%
  \tikzstyle{fancytitle} =[fill=white, text=black, rectangle, rounded corners, draw= black]%
  \noindent%
  \begin{tikzpicture}%
    \node [pmbox] (box){%
      \begin{minipage}{\textwidth}%
        \text{#1}%
      \end{minipage}%
    };%\
    \node[fancytitle, right = 10pt] at (box.north west) { \textbf{Problem}
    };%
    \end{tikzpicture}}

\def\theoremframecommand{\pmbox}
\newshadedtheorem{problem}{Problem}

\theoremclass{Proof}
\theoremstyle{empty}
\theorembodyfont{\normalfont}
 \newcommand{\pppmbox}[1]{
  \tikzstyle{pmbox} = [rectangle, rounded corners, draw=black, fill=cyan!5, inner sep=10pt, drop shadow={fill=black, opacity=.2}]%
  \tikzstyle{fancytitle} =[fill=white, text=black, rectangle, rounded corners, draw= black]%
  \noindent%
  \begin{tikzpicture}%
    \node [pmbox] (box){%
      \begin{minipage}{0.95\textwidth}%
        \text{#1}%
      \end{minipage}%
    };%\
    \node[fancytitle, right = 10pt] at (box.north west) { \textbf{Proof}
    };%
    \end{tikzpicture}}

\def\theoremframecommand{\pppmbox}
\newshadedtheorem{proof}{Proof}

\theoremclass{Example}
\theoremstyle{empty}
\theorembodyfont{\normalfont}
 \newcommand{\embox}[1]{
  \tikzstyle{pmbox} = [rectangle, rounded corners, draw=black, fill=Dandelion!5, inner sep=10pt, drop shadow={fill=black, opacity=.2}]%
  \tikzstyle{fancytitle} =[fill=white, text=black, rectangle, rounded corners, draw= black]%
  \noindent%
  \begin{tikzpicture}%
    \node [pmbox] (box){%
      \begin{minipage}{0.95\textwidth}%
        \text{#1}%
      \end{minipage}%
    };%\
    \node[fancytitle, right = 10pt] at (box.north west) { \textbf{Example}
    };%
    \end{tikzpicture}}

\def\theoremframecommand{\embox}
\newshadedtheorem{example}{Example}


\theoremclass{Definition}
\theoremstyle{empty}
\theorembodyfont{\normalfont}

 \newcommand{\dpmbox}[1]{
  \tikzstyle{pmbox} = [rectangle, rounded corners, draw=black, fill=gray!5, inner sep=10pt, drop shadow={fill=black, opacity=.2}]%
  \tikzstyle{fancytitle} =[fill=white, text=black, rectangle, rounded corners, draw= black]%
  \noindent%
  \begin{tikzpicture}%
    \node [pmbox] (box){%
      \begin{minipage}{0.95\textwidth}%
        \text{#1}%
      \end{minipage}%
    };%\
    \node[fancytitle, right = 10pt] at (box.north west) { \textbf{Definition}
    };%
    \end{tikzpicture}}

\def\theoremframecommand{\dpmbox}
\newshadedtheorem{definition}{Definition}
    
\theoremclass{Theorem}
\theoremstyle{empty}
\theorembodyfont{\normalfont}
\newcommand{\tpmbox}[1]{
\tikzstyle{tpmbox} = [rectangle, rounded corners, draw=black, fill=gray!5, inner sep=10pt, drop shadow={fill=black, opacity=.2}]%
\tikzstyle{fancytitle} =[fill=white, text=black, rectangle, rounded corners, draw= black]%
\noindent%
\begin{tikzpicture}%
\node [tpmbox] (box){%
  \begin{minipage}{0.95\textwidth}%
    \textit{#1}%
  \end{minipage}%
};%\
\node[fancytitle, right = 10pt] at (box.north west) { \textbf{Theorem}
};%
\end{tikzpicture}}

\def\theoremframecommand{\tpmbox}
\newshadedtheorem{theorem}{Theorem}

\theoremclass{Lemma}
\theoremstyle{empty}
\theorembodyfont{\normalfont}
\newcommand{\lmbox}[1]{
\tikzstyle{lmbox} = [rectangle, rounded corners, draw=black, fill=gray!5, inner sep=10pt, drop shadow={fill=black, opacity=.2}]%
\tikzstyle{fancytitle} =[fill=white, text=black, rectangle, rounded corners, draw= black]%
\noindent%
\begin{tikzpicture}%
\node [lmbox] (box){%
  \begin{minipage}{0.95\textwidth}%
    \textit{#1}%
  \end{minipage}%
};%\
\node[fancytitle, right = 10pt] at (box.north west) { \textbf{Lemma}
};%
\end{tikzpicture}}

\def\theoremframecommand{\lmbox}
\newshadedtheorem{lemma}{Theorem}


\theoremclass{Proposition}
\theoremstyle{empty}
\theorembodyfont{\normalfont}
\newcommand{\propmbox}[1]{
\tikzstyle{tpmbox} = [rectangle, rounded corners, draw=black, fill=orange!5, inner sep=10pt, drop shadow={fill=black, opacity=.2}]%
\tikzstyle{fancytitle} =[fill=white, text=black, rectangle, rounded corners, draw= black]%
\noindent%
\begin{tikzpicture}%
\node [tpmbox] (box){%
  \begin{minipage}{0.95\textwidth}%
    \textit{#1}%
  \end{minipage}%
};%\
\node[fancytitle, right = 10pt] at (box.north west) { \textbf{Proposition}
};%
\end{tikzpicture}}

\def\theoremframecommand{\propmbox}
\newshadedtheorem{proposition}{Proposition}

\theoremclass{Solution}
\theoremstyle{empty}
\theorembodyfont{\normalfont}

 \newcommand{\spmbox}[1]{
  \tikzstyle{spmbox} = [rectangle, rounded corners, fill=green!5, draw=black, inner sep=10pt]%
  \tikzstyle{fancytitle} =[fill=white, text=black, rectangle, rounded corners, draw= black]%
  \noindent%
  \begin{tikzpicture}%
    \node [spmbox] (box){%
      \begin{minipage}{0.95\textwidth}%
        \text{#1}%
      \end{minipage}%
    };%\
    \node[fancytitle, right=10pt] at (box.north west) {\textbf{Solution}%
    };%
    \end{tikzpicture}}
 
\def\theoremframecommand{\spmbox}
\newshadedtheorem{solution}{Solution}

% \begin{figure}
%     \centering%
%     \begin{tikzpicture}
%         \node at (0,0) (a) {a};
%         \node at (1,0) (b) {b};
%         \draw (a)--(b);

%         \node[above = \baselineskip of current bounding box] {%
%             \begin{minipage}{0.9\linewidth}
%                 Text or whatever you want to put here.

%                 \begin{equation*}
%                     1 = 2
%                 \end{equation*}
%             \end{minipage}%
%         };
%         \draw (current bounding box.north west) rectangle (current bounding box.south east);
%     \end{tikzpicture}
% \end{figure}


%%%%%%%%%%%%%%  PAGE SETUP %%%%%%%%%%%%%%%%%
% LaTeX has big default margins
\usepackage[margin=1.5in]{geometry}

% The following sets up some headers
\usepackage{fancyhdr}
\pagestyle{fancy}
\lhead{Algebra Notes} % Left Header
\rhead{\thepage} % Right Header
\cfoot{} % Center Foot (empty)

%%%%%%%%%%%%% SHORTCUTS %%%%%%%%%%%%%%%%%%%%
% You can define your own shortcuts too.
% Examples of custom commands
\newcommand{\half}{\frac{1}{2}}
\newcommand{\cbrt}[1]{\sqrt[3]{#1}}

% Document content begins here
\begin{document}
% \usegdlibrary{trees}



% Set up a title
\title{Combinatorics (Additive, extremal, graph, yayayayayay)}
\author{Raymond Xu}
\date{\today}
\maketitle

\begin{abstract}
  % This is an example of a \LaTeX\ document,
  % complete with theorems and headers and the like.
\end{abstract}

% This line makes a ToC
\tableofcontents

% This line starts a new page
\eject

Yeah so a lot of combinatorics is about finding valid constructions. Actually this is a trend in all of math, so a background in what a construction is is great. 

\section{Set Theory}

\section{Extremal Combinatorics}

\begin{definition}[Sunflower]
    A sunflower is a collection of sets $S_1, S_2, \ldots, S_k$ such that $S_i \cap S_j$ is constant over all $i,j$. 
\end{definition}

It's called a sunflower because it looks like one when you draw it: 

\begin{theorem}[Erdos-Rado Sunflower Lemma]
    Given a family $\mathcal F$ of size $s$ sets, we may find a sunflower of size $k$ if $|\mathcal F| > s!(k-1)k^{s-1}$
\end{theorem}

\begin{proof}
  Define $f(s,k)$ to be the maximal cardinality of $\mathcal F$ such that it does not contain a sunflower. Continue with induction on $s$, and suppose that no sunflower exists in $\mathcal F$. If $s=1$, clearly $f(1, k) \le k-1$. To induct $s-1 \to s$, let $A_1, \ldots, A_t$ be the maximum collection of pairwise disjoint sets in $\mathcal F$. Then define $A = \cup_{1 \le i \le t} A_i$ so that $|A| < sk$, or else we can select $k$ disjoint sets. Define $\mathcal F_a = \{S\not a : a \in S, S \in \mathcal F\}$. Bound $|\mathcal F| \le \sum_a |\mathcal F_a| \le f(s-1, k)\cdot sk \le s!(k-1)k^{s-1}$ as desired. 
\end{proof}

\begin{problem}
  
\end{problem}

% \begin{remark}
%   A better bound is $|\mathcal F| > s!(k-1)^s$, but asymptotically it doesn't matter. 
%   This bound is also very weak, though the lower bound is $(k-1)^s$ which is a difference of $s!$. 
% \end{remark}




\section{Standard Olympiad Combinatorics}

\begin{problem}[2025 IMO 6]
  Consider a $2025\times2025$ grid of unit squares. Matilda wishes to place on the grid some rectangular tiles, possibly of different sizes, such that each side of every tile lies on a grid line and every unit square is covered by at most one tile.

Determine the minimum number of tiles Matilda needs to place so that each row and each column of the grid has exactly one unit square that is not covered by any tile.
\end{problem}

\begin{solution}
Color all uncovered squares in the $n^2 \times n^2$ grid. Define a path to be a sequence of squares $s_1, \ldots, s_k$ such that $s_{i+1}$ has a larger $x$ and $y$ coordinate than $s_i$. We can always partition the grid into paths $\mathcal P_1, \ldots, \mathcal P_j$ such that every square is included in exactly $1$ path. 
\\
\textbf{Claim: }The minimum number of rectangles needed to satisfy the condition is exactly \[\sum_{i=1}^j |P_i| + 2(n-1) - 1 = n^2+2n-3\]
\\

\end{solution}

Let's generalize this.

\begin{problem}[Buffed to 3D]
  
\end{problem}

\subsection{Invariant}

\subsection{Bounding}

\section{Graphs}

\begin{lemma}[Lindström–Gessel–Viennot]
  Let $G$ be a weighted, acyclic graph, and let $A = \{a_1, a_2, \ldots, a_n\}$ and $B = \{b_1, \ldots, b_n\}$ be starting and ending vertices. Suppose that for any $i< j$ and $k<l$, every path $P_1 : a_i \to b_l$ intersects every path $P_2 : a_j \to b_k$. For two vertices, let $w(a,b)$ be the weight over the product of edges from $a \to b$. If \[M = \begin{pmatrix}
    w(a_1, b_1) & w(a_1, b_2) & \ldots & w(a_1, b_n) \\ w(a_2, b_1) & w(a_2, b_2) & \ldots & w(a_2, b_n) \\ \vdots & \vdots & \ddots & \vdots \\ w(a_n, b_1) & w(a_n, b_2) & \ldots & w(a_n, b_n)
  \end{pmatrix}\] Let $P_1, \ldots, P_n$ be an $n$-tuple of paths from $A$ to $B$, and let $\sigma$ be a permutation of $[n]$. Then \[\det(M) = \sum_{(P_1, \ldots, P_n) : A \to B} \text{sign}(\sigma(P)) \prod_{i=1}^n w(P_i)\] where $\text{sign}(\sigma) = (-1)^{N(\sigma)}$ where $N(\sigma)$ is the number of inversions. 
\end{lemma}

\begin{proof}
  By the definition of $M$, \[\det(M) = \sum_{\sigma \in S_n} \text{sign}(\sigma) \prod_{i=1}^n w(a_i, b_{\sigma(i)}) \] \[ =  \sum_{\sigma \in S_n} \text{sign}(\sigma) \sum \{w(P) : P \text{ is an $n$-tuple of paths }\}\] \[ =  \sum_{P = (P_1, \ldots, P_n) : A \to B} \text{sign}(\sigma(P)) w(P) \]
\end{proof}

\begin{problem}[2024 Putnam A6]
  Let $c_0,\,c_1,\,c_2,\,\ldots$ be a sequence defined so that
\[
\frac{1-3x-\sqrt{1-14x+9x^2}}{4}=\sum_{k=0}^\infty c_kx^k
\]for sufficiently small $x$. For a positive integer $n$, let $A$ be the $n$-by-$n$ matrix with $i,j$-entry $c_{i+j-1}$ for $i$ and $j$ in $\{1,\,\ldots,\,n\}$. Find the determinant of $A$.
\end{problem}

\begin{solution}[Solution 1]
  Weight every edge of the cartesian plane $(a,b) \to (a+1, b)$ with $5$ if $b \equiv a\pmod 2$ and $2$ if $b \equiv a+1 \pmod 2$. 
  \\
  \textbf{Claim: }The sum of the paths from $(0,0)$ to $(n-1, n-1)$ is exactly $c_n$. 
  \\ \textit{Proof. }Let $d_i$ be the sum of the weights with the $2$'s and $5$'s reversed. Then \[c_n = 5(c_{n-1}d_1 + c_{n-2}d_2 + \ldots + c_{1}d_{n-1})\] \[d_n = 2(c_{n-1}d_1 + \ldots + c_1d_{n-1})\] so $d_n = \frac{2}{5}c_n$ for $n\ge 2$, so \[c_n = 2(c_1c_{n-1} + \ldots + c_{n-1}c_1) + 3c_{n-1}\] 

  Let $a_i = (n-i+1, n-i+1)$ and $b_i = (n+i-1, n+i-1)$ so applying LGV gives $10^{n \choose 2}$ as the answer. 
\end{solution}

\begin{solution}[Solution 2]
  
\end{solution}

\section{Additive Combinatorics}






\end{document}
